\section{What cognitive restructuring asks of the patient}\label{sec:demands}
Consider the standard thought-record exercise. The patient must (a) \emph{notice} an automatic thought while it is occurring, (b) hold the thought, the supporting evidence, the contradicting evidence and a candidate alternative in mind at once, (c) suppress the pull of the automatic thought while evaluating it, and (d) shift from being inside the feeling to observing it. Mapped onto the best-known model of executive function, these are three separable components: updating and holding information (working memory), inhibiting a prepotent response, and shifting between mental sets \citep{miyake2000unity}.

\paragraph{Laboratory evidence that reappraisal draws on executive resources.} In event-related-potential studies, high working-memory load eliminated the effect of reappraisal on the late positive potential \citep{gan2017load,adamczyk2022load}. Working-memory capacity predicted reappraisal success \citep{schmeichel2008memory}, and reappraising late in an emotional response, once emotion is already intense, costs self-control resources \citep{sheppes2008costs}. People also choose differently as intensity rises: they prefer distraction to reappraisal in high-intensity negative situations and reappraisal in low-intensity ones \citep{sheppes2011choice}.

\paragraph{Evidence that it works, and can be trained, when load is low.} Instructed reappraisal reduces negative experience in the laboratory \citep{hofmann2009anxiety} and can be trained \citep{denny2014training}. So the claim is not that reappraisal fails; it is that it \emph{costs} executive resources, which matters when those resources are scarce.

\paragraph{Two clarifications.} First, this is a task analysis plus laboratory support, not a treatment finding. Clinicians scaffold the load (worksheets, a therapist who prompts for evidence, graded practice). Second, executive function is not a single resource; the argument concerns working memory, inhibition and shifting, not general intelligence. We return to both in Section~\ref{sec:objections}.
